\section{Five loop generations: arm L}\label{sec:arm-l}

Arm L iterates arm MU's conditioning step for five generations, feeding each scored proposal
back as the next parent. Iteration keeps the proposer off the anchor: 11 of 49 conditioned outputs sit on it,
against 11 of 17 at generation 0. Three of the four lineages end exactly on the family argmax
and none goes above it.

\textbf{Design.} Both prompts were hashed, and the stepper and scorer were committed
alongside. Generation 0 is the bare prompt, sampled
fresh inside this arm. Its prompt is byte-identical to the companion paper's unconditioned
arms, served on the same channel and wrapper. Generations 1 to 5 use arm MU's
registered mutation prompt verbatim, instantiated with the parent packing and its score. Two
discriminating cells, $N = 13$ and $N = 31$, are crossed with two archive regimes. GREEDY
retains only the single best valid proposal. DIVERSE retains up to five distinct-valued ones.
That gives four lineages of six rounds at population 5, which is 120 invocations. Two of the
120 were rejected by the runtime before reaching the model, and are counted as invalid rather
than dropped. The parent for population slot $i$ is
$\text{archive}[i \bmod |\text{archive}|]$, so the loop carries no sampling randomness of its
own. All four lineages are evaluable under the registered floor.

This is a minimal analogue. It has no islands, no crossover and no program representation.
Its only evaluator signal is the parent's own score.

\textbf{Registered predictions.} L1 predicted that the modal valid output at generation 0
equals the anchor in every lineage. L2 predicted that dissolution survives
iteration, so that the pooled on-prediction rate over generations 1 to 5 falls below the
generation-0 rate. L3 predicted that the pooled conditioned on-prediction rate is higher under
GREEDY than under DIVERSE, on the reading that loop diversity is scaffold-produced. L4 was
registered open with no direction predicted: report best-of-run against the family argmax.
The falsifier fires if the pooled conditioned rate meets or exceeds the generation-0 rate in
both cells.

\begin{table}[!ht]
\caption{Arm L by lineage. Valid outputs over invocations, and on-prediction counts, at
generation 0 and across the five conditioned generations, with each lineage's best scoring
valid output, the generation at which that best first appeared, and its cell's family
argmax. Pooled: 11 of 17 on-prediction at generation 0, 65\% (95\% Wilson 41--83\%), against
11 of 49 conditioned, 22\% (13--36\%).}
\label{tab:arm-l}
\begin{center}
\footnotesize
\setlength{\tabcolsep}{4pt}
\begin{tabular}{@{}lccccccc@{}}
\toprule
Lineage
& \begin{tabular}[b]{@{}c@{}}Gen-0 valid\\/ invocations\end{tabular}
& \begin{tabular}[b]{@{}c@{}}Gen-0\\on-prediction\end{tabular}
& \begin{tabular}[b]{@{}c@{}}Conditioned valid\\/ invocations\end{tabular}
& \begin{tabular}[b]{@{}c@{}}Conditioned\\on-prediction\end{tabular}
& \begin{tabular}[b]{@{}c@{}}Best\\of run\end{tabular}
& \begin{tabular}[b]{@{}c@{}}Best first\\reached (gen)\end{tabular}
& \begin{tabular}[b]{@{}c@{}}Family\\argmax\end{tabular} \\
\midrule
13-GREEDY & 5/5 & 2 & 4/25 & 0 & $1.7761424$ & 0 & $1.7761424$ \\
13-DIVERSE & 4/5 & 2 & 16/25 & 5 & $1.7761424$ & 5 & $1.7761424$ \\
31-GREEDY & 4/5 & 3 & 18/25 & 3 & $2.6104167$ & 5 & $2.7485281$ \\
31-DIVERSE & 4/5 & 4 & 11/25 & 3 & $2.7485281$ & 2 & $2.7485281$ \\
\midrule
Pooled & 17/20 & 11 & 49/100 & 11 & & & \\
\bottomrule
\end{tabular}
\end{center}
\end{table}

\textbf{Validity, and what every rate in this arm conditions on.} Validity drops under
iteration. The pooled row of Table~\ref{tab:arm-l} gives the two totals. The weakest lineage,
13-GREEDY, carries the loss. Conditioned rounds lose validity as well as accuracy, so every
rate here counts only survivors. The 11 of 49 on-prediction rate is one such rate. Over all
invocations it is 11 of 100. That is 6--19\%. At generation 0 it was 11 of 20. That is
34--74\%. The registration set no validity prediction, so the loss is a post hoc
observation.

\textbf{L2 confirmed, with unequal support.} Sum the two lineages at each cell of
Table~\ref{tab:arm-l}. At $N = 31$ the generation-0 on-prediction rate is 7 of 8. Conditioned,
it falls to 6 of 29. That is 21\% (10--38\%). At
$N = 13$ it starts at 4 of 9. Conditioned, it is 5 of 20. That is 25\% (11--47\%). Both drops
are point estimates. The falsifier did not fire. Once a lineage's parent sits
off the anchor, later rounds condition on a parent at the family argmax, arm MU's second
condition, so L2 is not independent of MU. Of the 100 conditioned rows, 34 carry a parent at
the family argmax: 25 in 13-GREEDY and 9 in 31-DIVERSE, none in the other two lineages.
13-GREEDY returned 4 of 25 valid under that
condition where arm MU's argmax-parent rows returned 26 of 45, a gap this paper does not
explain; 31-DIVERSE returned 7 of 9.

\textbf{L1 was not met.} It fails on one lineage. The 13-DIVERSE seed round in
Table~\ref{tab:arm-l} splits evenly, and the registered tie rule counts the tie against the
prediction. The other three lineages pass. The rule set four per-lineage bars rather than one
pooled bar per cell. Each bar ran over the four or five valid proposals that lineage's seed
round returned.

\textbf{L3 was not met.} Per cell, GREEDY and DIVERSE give 0 of 4 against 5 of 16 at
$N = 13$ and 3 of 18 against 3 of 11 at $N = 31$. Pool each archive regime over both cells.
The conditioned
on-prediction rate is 3 of 22 under GREEDY. That is 14\% (5--33\%). Under DIVERSE it is 8 of
27. That is 30\% (16--48\%). The difference runs opposite to the
registered direction and is not separable from chance. A post hoc Fisher exact test on the
pooled $2 \times 2$ table gives $p = 0.303$. The odds ratio is 0.375. The comparison is also
confounded: lineages differ in parent identity as well as archive width. 13-GREEDY
also contributes the fewest conditioned valid proposals of any lineage. The outcome is a
disconfirmation of L3, not evidence for its converse.

\textbf{L4: three of the four lineages reach the family argmax exactly, and none clears it.}
The best-of-run and family argmax columns of Table~\ref{tab:arm-l} carry the comparison. Both
lineages at $N = 13$ and 31-DIVERSE land exactly on their cell's family argmax, and 31-GREEDY
stops short of it. The best of 13-GREEDY is a generation-0 output, one family-argmax packing in its unconditioned seed round that
the loop never bettered. The other three lineages reach their best under conditioning. Of the 49 valid
conditioned outputs the table counts, 47 are valid at $10^{-9}$, and none of those exceeds
the best value the recipe family can express by more than $10^{-6}$. The largest excess among
the 49 is $1.3\times10^{-15}$, float noise, on five 31-DIVERSE rows. One post-sampling scorer
correction, affecting only the
$10^{-9}$ test and no verdict, is disclosed in Supplement~S2.
